% Sección 3
\section{Por qué en Manizales es más difícil}

\subsection{Una ciudad de ladera con pocas entradas}

Manizales está construida sobre una cuchilla, y casi todas sus vías de acceso suben o bajan por ladera. Esa descripción sale del conocimiento que el equipo tiene de su propia ciudad. No encontramos un estudio que mida cuánto depende Manizales de pocas vías de acceso, así que esa dependencia es por ahora inferencia nuestra.

Un análisis de la Universidad Javeriana de Cali sobre el sismo coincide con esa lectura. El estudio agrupó los lugares afectados en clústeres y puso a Manizales en uno que llama «urbano de ladera», con 361 puntos en el registro que usa y a 81 km del epicentro. Ese registro es abierto y no oficial.

El terreno también decide qué vehículo llega a cada sitio. Donde no entra un camión puede entrar un campero o una moto, y esa diferencia pesa en la última milla de los barrios de ladera y de las veredas. Por eso el prototipo trabaja con una flota mixta y con reglas de acceso por tipo de vehículo. Esta parte es inferencia nuestra y no viene de una medición.

\subsection{Cierres por derrumbe antes del sismo}

Los derrumbes cierran las vías de Manizales también sin terremoto. En marzo de 2017 un derrumbe cerró la vía Manizales–Fresno, que es la salida hacia Bogotá, en el sector Brisas, en el límite con Caldas. En abril de 2017, los deslizamientos en los barrios de ladera Aranjuez, Camilo Torres y Bajo Persia dejaron 17 muertos y 75 viviendas colapsadas, y obligaron a desalojar unas 400. La Cruz Roja estimó entonces que habría que reubicar unas 500 familias y revisar 11 barrios en riesgo.

Los cierres siguieron en los años siguientes. En marzo de 2021, un deslizamiento de unos 4.000~m³ por lluvias cerró por completo la doble calzada Manizales–Chinchiná en el kilómetro 11+500, y el tránsito se desvió por la vía antigua, en el sector de La Siria. En marzo de 2023, la vía hacia Mariquita se cerró de forma total e indefinida por lluvias y deslizamientos. Hubo que remover material hasta Las Letras, y el desvío fue por el Alto de La Línea. En junio de 2023, otro deslizamiento cerró la vía a Bogotá en el Alto de Letras. El INVÍAS recomendó entonces desviarse por Pereira, Armenia y La Línea, o por La Dorada y Medellín. En diciembre de 2024, un derrumbe cerca de Padua (Tolima) volvió a bloquear la vía hacia Fresno.

No todos los cierres son por derrumbe. En enero de 2018, el puente del Irra, en la vía hacia Medellín, se cerró de forma preventiva e indefinida por fallas estructurales, con desvíos por Supía, Riosucio y Anserma o por La Virginia y Pereira. Estos antecedentes son los que encontramos en prensa y en la ANI. La investigación no revisó los años 2011 y 2022, y no halló la duración de varios de estos cierres.

La ciudad tiene también capacidad propia frente al riesgo. En febrero de 2026, Manizales fue certificada como «Nodo de Resiliencia» dentro de la iniciativa MCR2030 de la UNDRR, la oficina de Naciones Unidas para la reducción del riesgo de desastres. Según Semana, la ciudad cuenta con un sistema de alerta multiamenaza y ha hecho más de 1.000 intervenciones de reducción del riesgo. La misma nota dice que no hubo muertes por deslizamientos en 2024 ni en 2025 hasta su fecha de publicación. El IDEA de la UNAL aporta además un indicador de alerta temprana, el A25, que usa la lluvia acumulada de 25 días con umbrales amarillo, naranja y rojo. Para la recuperación del territorio esta capacidad cuenta, porque un sistema de despacho puede apoyarse en ella en vez de partir de cero. Esa última idea es inferencia nuestra.

\subsection{La cadena de frío}

La cadena de frío depende de la energía. Al comienzo de la emergencia, la demanda de energía no atendida en Caldas, Quindío y Risaralda fue del 82~\%, y el 13 de agosto todavía era del 5,9~\%. El director ejecutivo de ABACO explicó que, sin energía, la red no podía manejar perecederos, y que por eso priorizó enlatados y alimentos de larga vida. La Cruz Roja pidió además no donar perecederos.

Nuestra lectura es que los perecederos quedaron casi fuera de la ayuda en los primeros días. Es inferencia nuestra, porque no encontramos ninguna cifra pública de alimentos donados que se dañaran o vencieran en Manizales. Como referencia nacional, una nota de Noticias RCN atribuye al Departamento Nacional de Planeación (DNP) una estimación de 9,76 millones de toneladas de alimentos perdidos o desperdiciados al año, el 34~\% de la oferta del país. Nadie del equipo ha abierto el documento del DNP.

\subsection{Cómo se organiza el territorio}

Manizales se divide en 12 comunas en la zona urbana y 7 corregimientos en la zona rural. Según el consolidado de juntas de acción comunal de la Alcaldía, la ciudad tiene 204 juntas, que corresponden a 131 barrios urbanos y 73 veredas. El documento no trae nombres ni coordenadas. La población proyectada para 2025 es de 472.087 habitantes, según una página que reproduce la proyección del DANE. Falta confirmar esa cifra en el propio DANE.

Las juntas de acción comunal son la forma en que barrios y veredas se organizan frente a la Alcaldía. Para DespachoX son la lista natural de quienes piden, porque la ficha del reto nombra a los comités barriales entre sus usuarios. Las juntas existen antes y después de una emergencia. Por eso un canal de pedidos montado sobre ellas podría seguir sirviendo cuando termine la fase de respuesta, que es lo que la mentoría llama recuperación del territorio. Esta idea es una hipótesis por validar con las propias juntas.

Las veredas pesan en este problema por dos razones documentadas. ABACO nombró la llegada a las zonas rurales como su mayor reto, y el Registro Único de Damnificados dejó por fuera a familias rurales dispersas. Un despacho que acepte pedidos colectivos por junta podría llegar a familias que no figuran en ese registro. Esa es una propuesta del equipo, tomada de las lecciones de otros desastres de la región.

La gente nombra los lugares a su manera. Un pedido puede decir «por la Siria» o «abajo de la Terminal», y quien despacha tiene que saber dónde queda eso. En Haití, en 2010, el conocimiento local de los nombres de lugar fue decisivo para ubicar los mensajes de ayuda. En Manizales todavía no sabemos qué comunas concentraron a los damnificados ni cuáles quedaron aisladas. Los únicos lugares que nombra la prensa son los cuatro puntos de censo de Milán, Chipre, la avenida Santander y el centro.
